\section{Equity, Accessibility, and AI Literacy}\label{sec:equity}

\subsection{Equity of access}
\begin{itemize}
  \item No student is required to purchase an AI tool or subscription. Where a course expects AI use (Levels~3--4), the faculty member ensures an accessible institutional or free option exists, or provides an alternative.
  \item The University will consider connectivity and device constraints, especially for students in low-bandwidth areas, by providing offline-friendly alternatives and on-campus access points. \todoverify{Resource feasibility.}
  \item Policies and detection practices must not disadvantage non-native English writers \citep{liang2023} or students using assistive technologies.
\end{itemize}

\subsection{Accessibility}
AI is encouraged as an accessibility aid (captioning, transcription, text-to-speech, simplification, translation). Use of AI for disability accommodation does not count as a violation of any AI-use level and is not required to be disclosed beyond what is needed to arrange the accommodation. Course materials produced with AI should be checked for accessibility.

\subsection{Bias and responsible interpretation}
Users are reminded that AI outputs can reflect bias, be inaccurate, or omit local and cultural context [CITATION NEEDED]. Course activities should include critical evaluation of outputs, including relevance to Philippine and Visayan contexts and languages.

\subsection{AI literacy}
\begin{itemize}
  \item \textbf{Students:} AI literacy is integrated into orientation and general-education or discipline courses: how AI works at a basic level, capabilities and limits, prompting, verification, ethics, privacy, and disclosure.
  \item \textbf{Faculty and staff:} see Section~\ref{sec:faculty} for development pathways.
  \item \textbf{Everyone:} a short, shared orientation and the quick-reference guides (Appendix~\ref{app:quickref}) provide a common baseline.
\end{itemize}

\subsection{Wellbeing}
Users are encouraged to be aware of overload and stress from rapid technology change; wellbeing supports are available through existing University services.
